\section{Scope and Systematization Method}\label{sec:methodology}\label{sec:background}
\subsection{Research questions, eligibility and threat polarity}
We organize the systematization around four questions: which hidden targets, threat polarities and auxiliary observables appear across LLM and agent systems (RQ1); how signal generation origin relates to receiver observation permissions (RQ2); which holdouts establish portability across tasks, configurations and hardware (RQ3); and when a component intervention protects an end-to-end workflow at a bounded utility and cost (RQ4).

Core inclusion requires five conjunctive criteria: (C1) LLM or agent execution, (C2) a hidden target property, (C3) an execution-dependent auxiliary observable rather than ordinary output semantics alone, (C4) specified observer permissions, and (C5) an empirical inference evaluation. Mixed-feature papers qualify only if they isolate an auxiliary-only C3 component; thus coding fingerprints, WildAuth, ZeroPrint, endpoint-content fingerprints and StepFinder remain adjacent boundary comparators \cite{ghaleb2026codingfingerprint,wang2026dialogue,shao2026zeroprint,leshin2026endpoint,zhu2026stepfinder}. Table~\ref{tab:boundary-rules} states the observable eligibility rules.
\begin{table}[t]\centering\footnotesize
\caption{Observable-based eligibility. A hidden target alone does not make a semantic attack a side channel.}\label{tab:boundary-rules}
\begin{tabularx}{\columnwidth}{@{}p{.39\columnwidth}Y@{}}\toprule
Observed input & Decision and required distinction\\\midrule
Length, timing, cache response & Core when measured execution dependence supports hidden-target inference.\\
Own-site actions & Core attribution with site permission; no automatic access to other sites or private tool logs.\\
Reply states the secret & Direct disclosure; adjacent, not auxiliary-channel evidence.\\
Reply style/content identifies a user & Adjacent unless a separately evaluated auxiliary mechanism is established. WildAuth stays a mixed-view comparator.\\
Output-filter membership oracle & Core only for an isolated accept/reject or differential-response oracle tied to filter execution; semantic content alone remains adjacent.\\
Privileged loop/fault logs & Separate diagnostic component when C1--C5 hold; does not increase confidentiality evidence.\\\bottomrule
\end{tabularx}\end{table}


We use four narrative groupings of objectives and observer roles, not mutually exclusive classes. An unauthorized identity inference can be both attribution and a confidentiality concern. Each record separately specifies the inference target, observer, authorization relation and protected interest:
\begin{enumerate}
\item \textbf{Type~I (Confidentiality and privacy protection):} An unauthorized observer attempts to infer private user prompts, personal traits, visited domains or proprietary model states; the defense protects user or tenant confidentiality (e.g., NetEcho, WRA, Early Bird, KVGov, CachePrune, OTRO \cite{netecho2025,jeong2026wra,earlybird2025,addagada2026kvgov,wu2026cacheprune,lee2026otro}).
\item \textbf{Type~II (Provenance and identity attribution):} An observer infers which model, agent harness or user identity generated a trace from timing, packet or DOM event structure (e.g., UI Traces, FP-Agent, AgentPrint, Lidar \cite{uitraces2026,wang2026fpagent,agentprint2025,wang2026lidar}).
\item \textbf{Type~III (Human--agent / bot-detection evasion):} The receiving platform operates a defensive classifier to detect automated bots, while the agent applies ``action humanization'' (e.g., AHB \cite{zhu2026turing}) to evade detection---inverting the defender/attacker polarity relative to Type~I privacy protection.
\item \textbf{Type~IV (Authorized operator diagnostics):} Privileged loop detection, serving root-cause localization or internal residual-stream monitoring (ClawGuard, cycle detection, LumiMAS, LLM RCA, MIRAGE \cite{gan2026clawguard,george2025cycles,solomon2026lumimas,scheinert2026rca,revankar2026mirage}), which use C1--C5 auxiliary signals for legitimate system telemetry rather than privacy exploitation.
\end{enumerate}

\subsection{Seed-anchored purposive corpus and exact reconciliation}
Because LLM and agent side-channel papers are scattered across security conferences, systems venues and preprints under heterogeneous terminology (``remote keylogging,'' ``prompt-cache leakage,'' ``workflow fingerprinting,'' ``agent humanization''), our corpus is constructed through \emph{purposive seed selection, venue-targeted proceedings screening (IEEE S\&P, USENIX Security, NDSS, ACM CCS, PMLR) and one-hop backward/forward citation snowballing} from 40 fixed seeds (2,237 backward reference occurrences and 24 distinct forward works screened; 38/40 forward seeds resolved; seed resolution and the number of distinct returned works are different counts), complemented by native database exports and bulk keyword searches (Appendix~\ref{app:search} and Supplement S1/S10).

To document why standalone Boolean keyword queries cannot replace seed-and-snowball discovery, we audited both our historical October 8--9 native IEEE/ACM exports and a completed October 21 supplementary search frame comprising four Semantic Scholar title/abstract bulk queries and two native arXiv searches (615 occurrences deduplicated to 565 unique records; Table~\ref{tab:corpus-reconciliation}). Those six supplementary keyword queries returned 17/24 of a fixed 24-paper core subset of our 40 seeds (\texttt{seed\_n=24}) and 25/60 of the October 8 snapshot. Because the 24 known papers are a core subset of the 40 initial seeds rather than an external holdout, 17/24 and 25/60 measure \emph{lexical overlap between narrow Boolean queries and our seed-anchored corpus}, not calibrated database recall: the two native arXiv queries required the exact hyphenated string \texttt{"side-channel"} (missing studies titled with ``prompt leakage,'' ``keylogging,'' ``timing attacks'' or ``fingerprinting''), and Semantic Scholar title/abstract matching did not search full texts of USENIX/NDSS PDFs.

\begin{table}[t]\centering\footnotesize
\caption{Item-by-item reconciliation of the 565-record October 21 search frame (before and after the 204-record exclusion audit) and the transition from the 60-family snapshot to all \FinalN{} coded families.}\label{tab:corpus-reconciliation}
\setlength{\tabcolsep}{3.5pt}
\begin{tabularx}{\columnwidth}{@{}Yrrr@{}}\toprule
Route / family category & Pre-audit & Audit $\Delta$ & Final (Indexed)\\\midrule
Excluded: no hidden C2/C3 target & 221 & $-4$ & 217 (0)\\
Excluded: outside LLM / diff.\ target & 175 & $0$ & 175 (0)\\
Existing snapshot core / adjacent & 23 / 1 & $+1 / 0$ & 24 / 1 (in 60)\\
New source-verified core families & 14 & $+1$ & 15 (\textbf{+15})\\
Adjacent source-reviewed contexts & 28 & $+2$ & 30 (\textbf{+7})\\
Adjacent metadata-only contexts & 28 & $0$ & 28 (0)\\
Background / proposal metadata & 58 & $0$ & 58 (0)\\
Missing abstract / unavailable text & 17 & $0$ & 17 (0)\\\midrule
\textbf{565-frame total / Net added families} & \textbf{565} & \textbf{0 net} & \textbf{565 (+22)}\\
\textbf{Coded collection}: Core / Adjacent (\textbf{Total}) & \multicolumn{3}{r}{Snapshot: 43 / 17 (\textbf{60}) $\to$ Full: \FinalCore{} / \FinalAdjacent{} (\textbf{\FinalN{}})}\\\bottomrule
\end{tabularx}
\par\smallskip\raggedright\footnotesize
The exclusion audit selected 177 high-risk pattern matches plus 27 reason-stratified records (seed 20261021), rather than all 396 initially excluded records. This enriched 204-record subset rerouted 4 records out of the 221 initial C2/C3 exclusions: (1) \texttt{SUP-0192} (MIRAGE \cite{revankar2026mirage}) moved to new core ($14+1=15$, all 15 added to the family index); (2) \texttt{SUP-0552} (InputSnatch arXiv alias) merged into existing family \texttt{inputsnatch2024} ($23+1=24$, $+0$ new families); and (3--4) \texttt{SUP-0191} and \texttt{SUP-0424} moved to source-reviewed adjacent contexts ($28+2=30$, of which 7 boundary-defining studies enter the adjacent index and 23 remain in the 565-record ledger). Thus $60 + 15\text{ (core)} + 7\text{ (adjacent)} = \FinalN{}$.
\end{table}

Table~\ref{tab:corpus-reconciliation} reconciles every record across the $60\to\FinalN{}$ transition. All \FinalCore{} core families (\FinalA{} workflow-origin, \FinalB{} inherited-runtime/application; 53 non-diagnostic and 5 Type~IV diagnostic) and \FinalAdjacent{} adjacent families are integrated into the unified corpus maps and intervention dispositions (Supplement S2/S5). Across the \FinalCore{} core families, 20 evaluate mitigations (16 snapshot families plus four additions). The result ledger is broader: 33 baseline records from 22 families, including six contextual comparators, plus six records from the four added mitigation families. Thus the 39-record ledger spans 26 families and 78 utility/cost field positions; the reporting-audit denominator is 20.

\subsection{Supplied coding records and agreement}
The supplied package records a primary AI-assisted coding workflow and a second coding column described by its provenance as a domain-expert review. We retain the recorded agreement summaries and supply their source columns. The files alone do not verify coder identity, blinding or independence from the primary decisions; we therefore report annotation agreement rather than newly verified independent human validation. The recorded scope is:
The saved screening sample has 93/100 three-class agreement ($\kappa=0.886$); the nine-item baseline audit has 134/144 cell agreement across 16 families ($\kappa=0.876$); and the 50-pair sample has 141/150 judgment agreement ($\kappa=0.916$). The latter covers all 13 same-family pairs plus 37 other within-cohort pairs. Full intervals and source locators remain in Supplement S8 and the annotation files.


The historical confidence intervals resample individual screening records, audit cells and pairwise judgments respectively. Cells within families and three judgments within a pair are dependent, so these intervals are conditional descriptive summaries, not cluster-robust generalization intervals. The 36 cells for four added families are outside the 144-cell agreement calculation. Agreement also does not establish coding correctness, and $\kappa$ depends on category prevalence.

\soktakeaway{SoK Takeaway 1 (Search and evidence scope)}{The supplementary queries returned 25/60 snapshot families and 17/24 known seed papers. These overlaps motivate complementary discovery methods; they neither estimate field-wide recall nor identify why each missing paper was absent. Objective, observer and authorization must be specified separately when comparing interventions.}
